\documentclass[header={Ancillary Note: Isometries},notheorems]{study-note}

\begin{document}

\StudyTitleBlock{N1}{Ancillary Note}{Isometries, isometric embeddings, isometric isomorphisms, and related concepts}

\vspace{0.8em}

\begin{focusbox}
This note is for quick memory support. It explains what an isometry is in the metric-space sense and in the linear-functional-analysis sense, and how this differs from being merely injective, surjective, bounded, or unitary.
\end{focusbox}

\tableofcontents

\section{The Basic Idea}

An isometry is a map that preserves distances exactly. In normed spaces this means it preserves the norm of differences, and for linear maps this reduces to preserving the norm of each vector.

\begin{defbox}{Isometry between metric spaces}
Let $(X,d_X)$ and $(Y,d_Y)$ be metric spaces. A map $T:X\to Y$ is an \textbf{isometry} if
\[
d_Y(Tx,Ty)=d_X(x,y)
\qquad \forall x,y\in X.
\]
\end{defbox}

\begin{warningbox}{Important distinction}
In a general metric-space setting, an isometry need not be linear. In functional analysis, when $X$ and $Y$ are normed spaces and $T$ is linear, one often says \textbf{linear isometry} to avoid confusion.
\end{warningbox}

\section{Linear Isometries In Normed Spaces}

\begin{defbox}{Linear isometry}
Let $X,Y$ be normed spaces and let $T\in\mathcal L(X,Y)$ be linear. Then $T$ is a \textbf{linear isometry} if
\[
\norm{Tx}_Y=\norm{x}_X
\qquad \forall x\in X.
\]
\end{defbox}

\begin{thmbox}{Equivalent formulation for linear maps}
If $T:X\to Y$ is linear, then the following are equivalent:
\begin{enumerate}
\item $T$ is an isometry of metric spaces for the metrics induced by the norms;
\item $\norm{Tx-Ty}_Y=\norm{x-y}_X$ for all $x,y\in X$;
\item $\norm{Tx}_Y=\norm{x}_X$ for all $x\in X$.
\end{enumerate}
\end{thmbox}

\paragraph{Reason.}
For linear maps,
\[
Tx-Ty=T(x-y),
\]
so preserving all distances is the same as preserving the norm of every vector.

\begin{thmbox}{Immediate consequences}
If $T:X\to Y$ is a linear isometry, then:
\begin{enumerate}
\item $T$ is injective;
\item $T$ is continuous and $\norm{T}=1$ unless $X=\{0\}$;
\item $T^{-1}:T(X)\to X$ is also an isometry.
\end{enumerate}
\end{thmbox}

\paragraph{Why injective?}
If $Tx=0$, then
\[
0=\norm{Tx}=\norm{x},
\]
so $x=0$.

\section{Isometric Embedding vs. Isometric Isomorphism}

\begin{defbox}{Isometric embedding}
A linear isometry $T:X\to Y$ is often called an \textbf{isometric embedding}. It identifies $X$ with the subspace $T(X)\subset Y$ without distorting norms.
\end{defbox}

\begin{defbox}{Isometric isomorphism}
If, in addition, $T$ is surjective, then $T$ is an \textbf{isometric isomorphism}. In that case $X$ and $Y$ are the same as normed spaces up to relabeling of points.
\end{defbox}

\begin{focusbox}
The slogan is:
\begin{itemize}
\item isometric embedding = norm-preserving copy of $X$ inside $Y$;
\item isometric isomorphism = $X$ and $Y$ are exactly the same geometry.
\end{itemize}
\end{focusbox}

\section{Related Concepts}

\begin{defbox}{Contraction and nonexpansive map}
A map $T:X\to Y$ between metric spaces is:
\begin{itemize}
\item \textbf{nonexpansive} if
\[
d_Y(Tx,Ty)\le d_X(x,y) \qquad \forall x,y\in X;
\]
\item a \textbf{contraction} if there exists $0\le q<1$ such that
\[
d_Y(Tx,Ty)\le q\,d_X(x,y) \qquad \forall x,y\in X.
\]
\end{itemize}
An isometry is the extreme case where the distance is preserved with equality.
\end{defbox}

\begin{defbox}{Bounded below}
A linear map $T\in\mathcal L(X,Y)$ is \textbf{bounded below} if there exists $c>0$ such that
\[
\norm{Tx}_Y\ge c\norm{x}_X \qquad \forall x\in X.
\]
Every linear isometry is bounded below with $c=1$, but the converse is false.
\end{defbox}

\begin{warningbox}{Do not confuse these notions}
Being injective is weaker than being an isometry.

Being surjective is weaker than being an isometry.

Being bijective and bounded is weaker than being an isometric isomorphism.
\end{warningbox}

\section{Standard Examples}

\begin{exbox}{Metric-space isometry that is not linear}
On $\R$ with the usual distance,
\[
T(x)=x+1
\]
is an isometry because
\[
\abs{T(x)-T(y)}=\abs{(x+1)-(y+1)}=\abs{x-y}.
\]
But it is not linear.
\end{exbox}

\begin{exbox}{Rotation on $\R^2$}
A rotation matrix
\[
R=\begin{pmatrix}
\cos \theta & -\sin \theta\\
\sin \theta & \cos \theta
\end{pmatrix}
\]
is a linear isometry of $\R^2$ with the Euclidean norm. It is also surjective, so it is an isometric isomorphism.
\end{exbox}

\begin{exbox}{Canonical embedding into the second dual}
For a normed space $X$, the canonical embedding
\[
J:X\to X^{**},
\qquad
(Jx)(\phi)=\phi(x)
\]
is a linear isometry. It need not be onto; it is onto exactly when $X$ is reflexive.
\end{exbox}

\begin{exbox}{Right shift on $\ell^2$}
Define
\[
S(x_1,x_2,x_3,\dots)=(0,x_1,x_2,x_3,\dots).
\]
Then
\[
\norm{Sx}_{\ell^2}=\norm{x}_{\ell^2},
\]
so $S$ is a linear isometry. But it is not surjective, because no vector in its range has a nonzero first coordinate. So $S$ is an isometric embedding, not an isometric isomorphism.
\end{exbox}

\begin{exbox}{Fourier-coefficient map on a separable Hilbert space}
If $H$ has a complete orthonormal system $(e_k)$, then
\[
T:H\to \ell^2,
\qquad
Tx=((x,e_k)_H)_{k\ge 1}
\]
is an isometric isomorphism. This is one concrete form of the statement that every separable Hilbert space is isometrically isomorphic to $\ell^2$.
\end{exbox}

\begin{exbox}{A bounded bijection that is not an isometry}
The map
\[
T:\R^2\to\R^2,
\qquad
T(x,y)=(2x,y)
\]
is linear, bounded, bijective, and has a bounded inverse, but it is not an isometry because it changes lengths.
\end{exbox}

\section{Hilbert-Space Terminology}

\begin{defbox}{Orthogonal and unitary operators}
On a real Hilbert space, a surjective linear isometry is often called \textbf{orthogonal}.

On a complex Hilbert space, a surjective linear isometry is called \textbf{unitary}.
\end{defbox}

\begin{thmbox}{Useful criterion in Hilbert spaces}
Let $H$ be a Hilbert space and let $U\in\mathcal L(H)$. If $U$ is linear, then the following are equivalent:
\begin{enumerate}
\item $U$ is unitary (or orthogonal in the real case);
\item $U$ is surjective and
\[
\norm{Ux}_H=\norm{x}_H \qquad \forall x\in H;
\]
\item
\[
(Ux,Uy)_H=(x,y)_H \qquad \forall x,y\in H,
\]
and $U$ is surjective;
\item
\[
U^*U=UU^*=I.
\]
\end{enumerate}
\end{thmbox}

\paragraph{Why inner products appear.}
In Hilbert spaces, preserving norms is equivalent to preserving inner products, via the polarization identity.

\section{Typical Confusions}

\begin{warningbox}{Common pitfalls}
\begin{enumerate}
\item A linear isometry need not be onto.
\item An injective bounded operator need not preserve norms.
\item A bijective bounded operator need not be an isometry.
\item In a metric-space context, isometry does not imply linearity.
\item In Hilbert spaces, ``unitary'' means more than bounded and invertible: it means exact preservation of the Hilbert geometry.
\end{enumerate}
\end{warningbox}

\section{What To Remember Quickly}

\begin{focusbox}
If you forget everything else, remember:
\begin{enumerate}
\item isometry = preserves distance exactly;
\item linear isometry = preserves norm exactly;
\item isometric embedding = norm-preserving copy inside a bigger space;
\item isometric isomorphism = onto isometry, so the two normed spaces are geometrically the same;
\item unitary/orthogonal operator = surjective linear isometry on a Hilbert space.
\end{enumerate}
\end{focusbox}

\end{document}
